\noindent\textbf{Por qué se usa este procedimiento.} La inercia respecto al eje horizontal es $I_z=\int_A y^2dA=b\int_{-h/2}^{h/2}y^2dy=bh^3/12$. La dimensión perpendicular al eje es la que aparece al cubo.

\begin{enumerate}

\item La altura es la dimensión perpendicular al eje, y entra al cubo.

\end{enumerate}

\noindent\textbf{Operaciones y resultados intermedios.}

\begin{enumerate}

\item Momento de inercia.

\[\fitmath{y_{1}=\frac{50\cdot {120}^{3}}{12}\approx 7200000\,\mathrm{mm^{4}}}\]

\end{enumerate}
